\documentclass{beamer}

\usepackage{xltxtra}
\usepackage{graphicx}
%\usepackage[utf8]{inputenc}
\usepackage{fontspec}
\setmainfont{Eurostile LT Pro}
\setsansfont{Eurostile LT Pro}
\usepackage{tikz}
\usepackage{comment}
\usetikzlibrary{arrows,decorations.pathmorphing,backgrounds,fit,positioning,automata,calc}

\colorlet{mygray}{black!40}
\colorlet{mylightgray}{black!20}
\colorlet{mygreen}{green!60!black}
\definecolor{metroblue}{RGB}{0,59,126}
\definecolor{metroyellow}{RGB}{249,174,0}
\newcommand{\insertauthorpic}{\includegraphics[width=1.446cm]{pics/philip.png}\includegraphics[width=1.446cm]{pics/thomas.png}}
\newcommand{\footertext}{%
  \fontsize{5}{5}\selectfont\begin{minipage}{3mm}\vspace*{-1pt}\raggedleft\insertframenumber\end{minipage}\hspace*{3mm}
  © METRO AG.\hspace*{3mm}\insertshortauthor{}: \insertshorttitle{}
}
\newcommand{\datesize}{\fontsize{10}{10}\selectfont}
\newcommand{\metroslide}{%
{\setbeamercolor{background canvas}{bg=metroblue}
\begin{frame}[plain]
\centering
\includegraphics{pics/metro2.png}
\end{frame}
}
}



%TODO
\mode<presentation>
{
\setbeamerfont{block title}{size={}}
\usefonttheme{professionalfonts} % using non standard fonts for beamer
\usefonttheme{serif} % default family is serif
\setbeamercovered{transparent}
\setbeamertemplate{navigation bars}{}
\setbeamertemplate{navigation symbols}{}
\setbeamersize{text margin left=4mm}
\setbeamersize{text margin right=4mm}
\setbeamertemplate{itemize items}[ball]
\setbeamertemplate{enumerate items}[ball]
\setbeamertemplate{headline}{}
\setbeamertemplate{footline}{
  \begin{tikzpicture}[remember picture,overlay]
  \draw[metroblue] ($(current page.south west)+(0.6,0.6)$) -- ($(current page.south east)+(-3.5,0.6)$);
  \node (logosmall) at ($(current page.south east)+(-0.2,-0.1)$) [anchor=south east] {\includegraphics[height=8mm]{metro.png}};
  \node[metroblue] (textline) at ($(current page.south west)+(0.8,0.13)$) [anchor=south west] {\footertext{}};
  \end{tikzpicture}
}
\defbeamertemplate*{title page}{metro}[1][]{
  \begin{tikzpicture}[remember picture,overlay]
  \node (logo) at (current page.north west) [anchor=north west] {\includegraphics[height=\paperheight]{left.png}};
  \node[white] (title) at ($(logo)+(-0.7,0.65)$) [anchor=south west] {\fontsize{20}{20}\selectfont\textbf{\inserttitle}\par};
  \node[white] (subtitle) [below=0.7 of title.west,anchor=west] {\fontsize{12}{12}\selectfont\textbf{\insertsubtitle}\par};
  \node[white] (date) [below=0.9 of subtitle.west,anchor=west] {\datesize\insertdate,};
  \node[white] (place) [below=0.42 of date.west,anchor=west] {\datesize\selectfont\location};
  \node[white] (authorpic) [below=0.1 of place.south west,anchor=north west] {\insertauthorpic};
  \node[white] (author) [right=0.1 of authorpic.south east,anchor=south west] {\datesize\selectfont\insertauthor};
  \node[white] (textline) at ($(current page.south west)+(0.8,0.13)$) [anchor=south west] {\footertext{}};
  \end{tikzpicture}
}
\defbeamertemplate*{frametitle}{metro}[1][]{
  \vskip2mm%
  \hspace*{3mm}\fontsize{16}{12}\selectfont\textbf{\textcolor{metroblue}{\insertframetitle}}
}
}


\begin{document}

\title[React and Reflux]{REACT AND REFLUX:}
\subtitle{Components, Actions, Stores, and State}
\author[P.\ Baues, T.\ Str\"oder]{Philip Baues and Thomas Str\"oder}
\date{28./29.\ September 2017}
\newcommand{\location}{\texttt{code.talks} Hamburg}

{\setbeamercolor{background canvas}{bg=metroblue}
\begin{frame}[plain]
\centering
\includegraphics{metro2.png}
\end{frame}
}

{\setbeamercolor{background canvas}{bg=metroblue}
\begin{frame}[plain]
\titlepage
\end{frame}
}

\begin{comment}
Abstract:
React is currently one of the most popular ingredients of modern web applications. It is a component-based JavaScript library supporting the decomposition of big web applications into small parts, especially following principles like single responsibility and separation of concerns. However, on its own it lacks concise support for transversal state handling when parts of an application should alter the state of another part and these parts are not in a father-child-like relationship. To this end, the Reflux JavaScript library adds actions and stores for state handling.

In this talk, we will have a look at general patterns how to use React components together with Reflux stores and actions to build a responsive web application. Here, the focus is on state: When and where to (not) have it and how to deal with it using the facilities React and Reflux are offering. We illustrate these concepts with a small web shop example. In addition to that, we also address common mistakes and how to avoid them.
\end{comment}

\begin{frame}\frametitle{Frame Title}
\centering
Content
\end{frame}

\begin{frame}\frametitle{Frame Title}
\centering
Content
\end{frame}

\begin{frame}\frametitle{Frame Title}
\centering
Content
\end{frame}

\begin{frame}\frametitle{Frame Title}
\centering
Content
\end{frame}

\begin{frame}\frametitle{Frame Title}
\centering
Content
\end{frame}

\begin{frame}\frametitle{Frame Title}
\centering
Content
\end{frame}

\begin{frame}\frametitle{Frame Title}
\centering
Content
\end{frame}

\begin{frame}\frametitle{Frame Title}
\centering
Content
\end{frame}

\begin{frame}\frametitle{Frame Title}
\centering
Content
\end{frame}

\begin{frame}\frametitle{Frame Title}
\centering
Content
\end{frame}

\begin{frame}\frametitle{Frame Title}
\centering
Content
\end{frame}

\begin{frame}\frametitle{Frame Title}
\centering
Content
\end{frame}

\end{document}

